% R15 component; only input paths and enumerated location/grammar prose are edited.
% Source: revisions/2026-10-04-r15/manuscript/introduction.tex; commit 1cb3cc9efe135966ec228dfaf51c6841a6dfc97c.
\section{Introduction}\label{sec:introduction}
Many economic decisions depend on continuation values that must themselves be computed. A household chooses consumption and portfolio risk while investing in its preferences. Recursive utility makes future utility part of the evaluation equation. A sophisticated consumer anticipates the choices of future selves, and a firm invests against the policies of its rivals. In each case, a proposed change in behavior is useful only if the continuation object is evaluated accurately enough for the relevant economic decision. The preference-formation problem of \citet{qi2025} brings these internal and external choices together.

This paper studies Neural Bellman Operators (NBO), a numerical evaluation and improvement map for these problems. A critic evaluates the continuation value of a specified policy. A feasible actor changes the actions entering the economic generator, holding that value and its derivative channels fixed. An independent error account evaluates the resulting policy. The numerical object includes its boundary conditions, utility domain, approximation scheme, information structure, and deployment rule. These restrictions are part of the method: they determine what the computed policy and its reported accuracy mean.

The economic role of the critic is the central question. Learning a continuation value can pool information across states and actions, but the fitting cost and approximation bias may exceed the benefit of that reuse. A positive policy gain does not establish that its critic was useful. Nor does an accurate derivative at training states establish an improvement along the distribution generated by the deployed candidate. We connect evaluation error to feasible action quality and then to economic payoffs, keeping the distribution, the evaluation target, and the numerical transfer explicit. The computational comparisons examine the same chain through direct alternatives to the learned evaluation block.

The analysis has three parts. First, the NBO specification separates policy evaluation from action improvement and states the boundary and utility conditions needed by both. Differential and positive-weight Bellman formulations give distinct numerical realizations of this map. Their error accounts distinguish policy-evaluation defects from incomplete action maximization; the monotone formulation supplies a viscosity-selection argument under its stated approximation hypotheses. These results concern the output errors of a numerical method, rather than convergence of a particular neural optimizer.

Second, a nonlinear capital economy permits a detailed connection between the computed controller and the original continuous-time problem. Log felicity, bounded consumption shares, additive diffusion, and bounded coupled production yield an analytical comparison policy and explicit propagation bounds. A policy-specific payoff identity, diffusion transfer, and simultaneous common-path inference measure the gain of a fitted candidate and its difference from another candidate. The absolute comparison with the optimum retains its own anchor gap. An exact finite Bellman identity places continuation error along the feasible change from the reference action to the candidate action, integrated under the candidate's occupation law. A computable maximization-gap bound records unfinished actor optimization even when global concavity of the learned continuation is unavailable. The resulting gain bound transfers payoffs to the diffusion without equating discrete and continuous costates.

Third, the computational comparison measures economic performance and the work used to attain a stated target. It executes all sixteen members of a declared finite distribution of complete algorithm streams for NBO, cached antithetic Raw improvement, direct policy optimization, and direct neural HJB fitting in dimensions ten and fifty. The experiment uses the same nonlinear capital model at higher volatility, with fixed architectures, training schedules, and stopping rules. All 128 executions return fitted policies. NBO's schedule-relative method lower gains exceed $0.000943$ and $0.000957$ in the two dimensions.

The direct payoff comparison establishes a material advantage over the specified neural HJB procedure. An independently proved common-innovation transfer sharpens the deterministic allowance on the original observations, retaining the original confidence family, training, and stopping decisions. The resulting NBO-minus-HJB intervals are contained in $[0.000179,0.000411]$ and $[0.001347,0.001650]$; both exceed the prespecified $10^{-4}$ economic margin. The original intervals are reported alongside this refinement. The four NBO--Raw/DPO contrasts remain unresolved, rather than establishing equality. In dimension fifty, every NBO, Raw, and DPO stream attains the online gain target $0.0005$. Their mean complete process work is approximately $357.80$, $346.22$, and $596.94$ seconds; HJB attains the target in no stream. NBO therefore improves on this direct residual-fitting implementation and uses less measured work than DPO for the declared target, while Raw remains cheaper. These findings concern the executed finite method distributions, not all possible training designs.

The mechanism assessment separately computes every term of the finite Bellman account on the returned NBO policies. It uses the candidate's occupation law and independent action-bridge continuation banks, including the training-to-assessment grid discrepancy. Its protected lower gains are approximately $-0.006778$ and $-0.004800$, so this sufficient account does not certify a positive mechanism gain at the executed precision. The critic-risk comparison with Raw is also inconclusive in both dimensions. This distinction keeps the method-payoff finding separate from a claim that costate regression has established a general risk reduction. The earlier capital comparisons, including their unfavorable raw-costate diagnostics and inconclusive direct payoff intervals, remain visible.

Absolute economic accuracy has its own account. The new NBO full-adapted-class regret upper bounds are at most $0.025971$ and $0.036843$, which do not certify proximity to the optimum at the $10^{-4}$ margin. An independent scalar Howard calculation instead supplies finite-grid residual and action-restriction bounds, while its deployed classical policies outperform the frozen historical neural policies on the measured scalar comparisons. The scalar NBO--DPO ranking reverses at one initial state. These observations identify the scope of the approximation evidence without changing the economic problem or treating a one-dimensional grid as a high-dimensional reference.

The information structure also has a concrete economic implication. A manager can observe capital at finite dates, apply a saved feasible consumption rule, and hold that choice until the next observation. A separate experiment verifies this measured-state implementation against its corresponding protected parent. All four parent comparisons have positive lower gains, and six of twelve measurement specifications retain a positive continuous-economy lower gain after their sensing allowances. At 4,096 decision dates with exact measurements, the lower gains are approximately 0.001051 and 0.001054 in dimensions 10 and 50. Measurement error can exhaust this margin. The result identifies the precision required by the specified controller rather than assuming that finite observations reveal latent innovations.

The original applications remain essential to the economic scope of NBO. Recursive utility changes the fixed-policy evaluation problem. Endogenous preference adjustment links an internal shadow value to consumption and portfolio demand. Temporal selves require a spike-deviation equilibrium against two continuation values, while strategic firms require a complete unilateral best response against fixed rivals. These applications share the evaluation--improvement structure but have different economic error criteria. Their full models, derivations, proofs, and numerical records are retained in the current appendices and the indexed reproduction record.
